\documentclass[11pt,a4paper]{article}
\usepackage[utf8]{inputenc}
\usepackage[T1]{fontenc}
\usepackage{lmodern}
\usepackage[margin=2.2cm]{geometry}
\usepackage{amsmath,amssymb,amsthm}
\usepackage{booktabs}
\usepackage{array}
\usepackage{enumitem}
\usepackage[hidelinks]{hyperref}

\setlist{itemsep=2pt,topsep=3pt}
\setlength{\parindent}{0pt}
\setlength{\parskip}{6pt}
\raggedright

\newtheorem{proposition}{Proposition}
\newtheorem{corollary}{Corollary}

\newcommand{\strat}[1]{\texttt{#1}}
\newcommand{\Ddecl}{D^{\mathrm{decl}}_t}
\newcommand{\Dobs}{D^{\mathrm{obs}}_t}
\newcommand{\Trep}{T^{\mathrm{rep}}(\Delta)}
\newcommand{\Tdiv}{T^{\mathrm{div}}(\Delta)}

\title{Replay first: which evals to run after a change,\\ in what order, and when to stop}
\date{}

\begin{document}
\maketitle
\vspace{-1.5cm}

\section{The problem}

An LLM agent system is made of agents (a model with a prompt, skills and tools), the harness code that runs them, and
configuration. A \emph{test}, or eval, sends the system a request and checks what it did: the final state, the tool
calls and the answer. Every test calls the model, so every test costs tokens.

A \emph{change} $\Delta$ turns the current version $V_0$ into a new version $V_1$: an edited prompt or skill, a tool
that works differently, a new setting. Because a model produces the outcome, we run each test $K$ times on $V_0$ and
call $f_t$ the share of those runs that fail. A test with $f_t \ge \frac12$ is one the model already fails, and it
cannot tell us anything about a change. A test \emph{reveals} $\Delta$ when it passes on $V_0$ and fails on $V_1$;
$\mathrm{Rev}(\Delta)$ is the set of such tests, and $\Delta$ is a \emph{regression} when that set is not empty.

The question is which tests to run after a change, in what order, and when to stop, so that we catch regressions
while paying as few tokens as possible. Our answer has three steps:
\begin{itemize}
  \item \textbf{Replay.} We record the model's answers on $V_0$ once. After a change, we rerun every test with those
        recorded answers in place of the model. A test that never asks the model anything new gets its outcome for
        free. A test that does ask something new is one the change reaches, and only those need the model.
  \item \textbf{Order.} A strategy puts the reached tests in an order, and the run stops at the first failure.
  \item \textbf{Stop.} A stopping rule may call the change clean before every reached test has run, with a
        threshold set on past changes so that the share of missed regressions stays under a target we choose.
\end{itemize}

\section{How the replay works}
\label{sec:replay}

\paragraph{A run, written down.} A run of test $t$ on version $V$ alternates calls to the model $M$ with code. By
\emph{code} we mean everything that is not the model: the harness, the tools and the services behind them, the
starting data, and the test's checks. Each call sends a \emph{request} $q_i$ and gets an \emph{answer} $a_i$:
\[
  q_1 = g^t_V(), \qquad a_i = M(q_i), \qquad q_{i+1} = g^t_V(a_1, \dots, a_i), \qquad X_t(V) = h^t_V(a_1, \dots, a_n).
\]
\begin{itemize}
  \item A request is everything the model sees in that call: the system prompt, the messages so far including every
        tool result, the tools on offer, and the decoding settings. Calls to a judge model that grades the answer
        are requests too.
  \item $g^t_V$ runs the code of $V$ up to the next call. $h^t_V$ runs it to the end and applies the test's checks;
        $X_t(V)$ is the outcome, pass or fail.
\end{itemize}

\paragraph{The cassette.} We run the suite on $V_0$ and keep every answer under the request that produced it. $Q_0$
is the set of those requests and $A_0(q)$ the answer stored for $q$; when repeated runs disagree, we keep the most
frequent answer. We call this store the \emph{cassette}. A request is looked up by a hash of all of it, so two
requests match only if they are identical. Tool-call ids that a server makes up are replaced by their order of
appearance before hashing, so that they do not break the match.

\paragraph{Replaying a test.} To replay test $t$ on $V_1$, we run it with $A_0$ in place of $M$. The code is the
code of $V_1$ and runs as it is; only the model's answers come from the cassette. One of two things happens:
\begin{itemize}
  \item every request is in $Q_0$. The run finishes without the model, with outcome $\hat{X}_t(V_1)$, and $t$
        \emph{replays to the end}. These tests form $\Trep$;
  \item the run reaches a request $q \notin Q_0$, its \emph{divergence point}. The cassette has no answer, the
        replay of $t$ stops there, and $t$ \emph{diverges}. These tests form $\Tdiv$, the tests the change
        \emph{reaches}.
\end{itemize}
The replay makes no model call, so it costs no tokens and takes only the time the code needs.

\paragraph{Where a change shows.} Since $g$ and $h$ are the code of $V_1$, a change shows wherever it alters what
the model sees or what the code does with the answers:
\begin{itemize}
  \item a change to what an agent gets up front (its prompt, its skills, its tools and their descriptions, its task)
        alters the first request of that agent's conversation, and the test diverges there;
  \item a change to a tool's code, a service, the starting data or a setting the code reads alters a tool result
        whenever the test's call hits the changed behaviour. The next request is then new, and the test diverges
        there. Tools are never replayed from the recording; they always run the new code;
  \item a change that only alters state the model never reads is caught by the checks in $h$: the test replays to
        the end and fails;
  \item a change that alters none of this leaves the run as it was on $V_0$, and the test replays to its recorded
        outcome.
\end{itemize}

\paragraph{Not prompt caching.} Model providers and servers offer \emph{prompt caching}, also called prefix caching:
they keep the computation for a prompt prefix they have seen before, and still generate a new answer. That makes a
call cheaper and faster, but every call is still made. The cassette is different: a request it holds makes no call
at all.

\section{When the replayed outcome is the real one}
\label{sec:conditions}

The replay stands in for the model, so it tells us about a real run only if two conditions hold.
\begin{itemize}
  \item \textbf{Condition R.} With the model served the way we serve it, the same request always gets the same
        answer: $M(q) = A_0(q)$ for every $q \in Q_0$.
  \item \textbf{Condition D.} Given the answers, the code is deterministic: $g^t_V$ and $h^t_V$ read no clock, draw
        no random numbers, and call no outside service whose reply can change.
\end{itemize}

\begin{proposition}
\label{prop:replay}
Under R and D, a test in $\Trep$, run on $V_1$ with the real model, sends the requests of its replay and gets the
same answers, so $X_t(V_1) = \hat{X}_t(V_1)$. A test in $\Tdiv$ matches its replay up to its divergence point.
\end{proposition}

\begin{proof}
By induction on $i$. The first requests match, because both runs execute the code of $V_1$, which is deterministic
by D. If the $i$-th requests match and are in $Q_0$, the model answers $M(q_i) = A_0(q_i)$ by R, so the $i$-th answers
match, and by D so do the next requests. A test in $\Trep$ ends with the same answers, and $h^t_{V_1}$ turns them into
one outcome. A test in $\Tdiv$ matches up to its first request outside $Q_0$.
\end{proof}

\begin{corollary}
\label{cor:safe}
Under R and D, every test that reveals the change either diverges or fails in the replay:
$\mathrm{Rev}(\Delta) \subseteq \Tdiv \cup \{t \in \Trep : \hat{X}_t(V_1) = \text{fail}\}$.
\end{corollary}

\begin{proof}
A test in $\Trep$ that reveals the change fails on $V_1$, and by the proposition its replay has the same outcome.
\end{proof}

Neither condition is assumed. We check both:
\begin{itemize}
  \item \textbf{D} belongs to the system under test: services simulated or recorded, no clock, no randomness, keys
        derived from the operation. It is checked for free: replayed against its own cassette, $V_0$ must reach the
        end of every test with the recorded outcome. A system whose tools call outside services meets D only if
        those services are recorded or simulated too.
  \item \textbf{R} belongs to the model and the way it is served. Even with deterministic decoding (temperature 0,
        always the most likely token), a GPU server is not exact to the bit: the batch a request shares with other
        requests, and the server's prefix cache, change the floating-point arithmetic and now and then a token. We
        measure R with the \emph{empty change}: $V_0$ run again exactly the way a changed version is run, compared
        with the cassette request by request, run by run, and outcome by outcome. Replay first needs only the last
        of the three.
  \item When exact answers matter, serve the model without its prefix cache and with one test at a time per server,
        and confirm it with the empty change.
  \item A cassette is valid only for the model and the serving setup it was recorded with. A hosted model can
        change without notice; running the empty change from time to time tells us when to record again.
\end{itemize}

\section{What the replay gives}
\label{sec:gives}

Under R and D:

\begin{center}
\begin{tabular}{p{7.2cm}p{8.4cm}}
\toprule
the replay finds & so \\
\midrule
$t \in \Trep$ passes & $t$ passes on $V_1$; no need to run it \\
$t \in \Trep$ fails, and $f_t < \tfrac12$ & regression, before any model call \\
$\Tdiv$ is empty, and no such failure & clean: no test can tell $V_1$ from $V_0$ \\
$t \in \Tdiv$ & $t$ needs the model, for its requests outside $Q_0$ \\
every $t \in \Tdiv$ with $f_t < \tfrac12$ is run and passes & clean, with nothing the change reaches left untested
  \\
\bottomrule
\end{tabular}
\end{center}

\paragraph{What a reached test costs.} A reached test runs with the model behind the cassette: any request already in
$Q_0$ is answered from the cassette, and only new requests go to the model. Its cost is
\[
  c_\Delta(t) = \sum_{q \,\in\, \mathrm{run}(t, V_1),\ q \notin Q_0} \mathrm{tokens}(q).
\]
Before any model call, the replay bounds it from above by the tokens of $t$ after its divergence point. Strategies
plan with that bound.

\paragraph{Compared with a plain answer cache.} A cache of answers in front of the model would give the same saving
to any rerun of the suite. What the replay adds is knowing, before any call, which tests the change reaches. That
buys three more savings:
\begin{itemize}
  \item tests the model already fails on $V_0$ are not run;
  \item the run stops at the first failing test, so the order of the reached tests decides what a regression costs;
  \item a stopping rule decides what a change that breaks nothing costs.
\end{itemize}
The divergence point also says where the change reaches the model: at the start of an agent's conversation, when
what the agent gets has changed, or after a tool result, when what a tool returns has changed.

\section{Ordering the reached tests}
\label{sec:order}

This is classic test prioritization on a smaller suite. Every strategy follows the same two rules:
\begin{itemize}
  \item a test with $f_t \ge \frac12$ is not run;
  \item the run stops at the first failure, since a test that passed on $V_0$ and fails on $V_1$ means a regression.
\end{itemize}
With the replay, only tests in $\Tdiv$ can be run, and a test costs $c_\Delta(t)$.

The strategies differ in what they read. A \emph{component} is a part of the system a change can touch: a function
or line of code, a tool and its description, a skill's procedure or text, a field of an agent, a setting, the
starting data. $C(\Delta)$ is the set of components the change touches. For each test we know:
\begin{itemize}
  \item $c(t)$, its cost in tokens on $V_0$, and $\rho_t$, a risk weight its author gave it;
  \item $\Ddecl$, the components its author declared it depends on;
  \item $\Dobs$, the components its runs on $V_0$ actually exercised, from traces and code coverage, and
        $\phi_t(d)$, the share of its runs that exercised component $d$.
\end{itemize}
A component matches a test by one rule used everywhere: a changed function or line matches the tests that ran it, a
skill the tests that loaded it, a tool the tests it was offered to, a setting the tests that read it.

\begin{itemize}
  \item \strat{random}: a random permutation, as a reference.
  \item \strat{declared}, classic selection on declared dependencies: take $S = \{t : \Ddecl \cap C(\Delta) \neq
        \emptyset\}$ and $a_t = |\Ddecl \cap C(\Delta)|$, sorted by $\rho_t$, then $a_t$, both descending, then by
        $c(t)$ ascending.
  \item \strat{observed}, the same on observed dependencies: take $S = \{t : \Dobs \cap C(\Delta) \neq \emptyset\}$,
        with $a_t = \sum_{c \in C(\Delta) \cap \Dobs} \phi_t(c)$, sorted by $a_t\, \rho_t / c(t)$, descending. A
        variant, \strat{observed\_diverse}, takes the tests of $S$ from different behaviour groups in turn.
  \item \strat{cost\_aware} sorts the whole suite by $\rho_t\, (\beta_0 + a_t + e_t)\, (1 + f_t) / c(t)$, where
        $e_t = 1$ if $t$ ran a changed line and $\beta_0$ is a small constant.
  \item \strat{ir\_tfidf} sorts by text similarity. TF-IDF turns a text into word weights, where a word counts more
        when it is frequent in this text and rare across the tests. The score is the cosine between the weights of
        the diff and those of the test's text (name, request, expected outcome).
  \item \strat{history\_ml} learns from past changes. A logistic regression, a linear model that outputs a
        probability, is fitted on pairs (past change, test) labelled ``this test revealed this change'', with
        features such as the change's static impact on the test by kind of component, and the test's cost, risk and
        shape. With $p_t$ its probability, tests are sorted by $p_t\, \rho_t / c(t)$.
  \item \strat{adaptive\_bayes}, our method, recomputes the order after every test:
  \begin{itemize}
    \item a hypothesis $h \in \{0,1\}^{C(\Delta)}$ says which changed components are broken. Each component starts
          with a prior probability $\pi_c$ that depends on its kind: higher for code, skill procedures, settings and
          agent fields, lower for prompt wording, skill prose and tool descriptions;
    \item a test passes only if it would pass on $V_0$ and no broken component makes it fail:
          $P(t \text{ fails} \mid h) = 1 - (1 - f_t) \prod_{c:\, h_c = 1} (1 - s_{tc})$;
    \item the \emph{sensitivity} $s_{tc}$, the chance that $t$ fails when $c$ is broken, is learned from past
          changes: $s_{tc} = \mathrm{clip}(r_k\, \ell_{tk}\, \phi_t(c)\, \tau_t,\ s_{\min},\ s_{\max})$, where $k$ is
          the kind of $c$, $r_k$ is how often tests that touched past changes of kind $k$ revealed them, $\ell_{tk}$
          is how much better than average $t$ did on that kind, and $\tau_t$ is a bonus when the test's text is close
          to the diff or its author declared the dependency;
    \item after each result, Bayes' rule gives the \emph{posterior}, the probability of each hypothesis given the
          results so far. Its \emph{entropy} $H(h)$ measures how unsure we still are about what is broken. The next
          test is the one whose result is expected to lower that uncertainty the most per token, tilted by risk:
          $\arg\max_t \frac{H(h) - \mathbb{E}[H(h \mid X_t)]}{c(t)}\, (1 + \gamma(\rho_t - 1))$;
    \item after a test passes, the sensitivity of near-duplicates of that test is multiplied by $\delta < 1$, since
          their passing would add little.
  \end{itemize}
  \item \strat{casper\_slices}, adapted from CASPER: it groups the tests with k-means (clusters of tests whose
        behaviour on $V_0$ looks alike: turns, tool calls, skills loaded, pass rate), runs a few representatives of
        each group on $V_1$, and then runs the groups whose failure rate moved the most.
  \item \strat{risk\_aware\_ts}, adapted from Risk-Aware Adaptive Evaluation: a bandit with Thompson sampling. At each
        step, a logistic regression on the results so far gives each test an estimate $\mu_t$; we draw a plausible
        failure probability $\theta_t \sim \mathrm{Beta}(1 + \alpha\mu_t + k_t,\ 1 + \alpha(1 - \mu_t) + n_t - k_t)$,
        where $n_t$ and $k_t$ count the runs and failures seen for $t$, and run the test with the largest
        $\rho_t\, \theta_t$.
  \item two orders that need nothing but the replay: \strat{cheapest}, by the cost bound, ascending, and
        \strat{late\_divergence}, by how late the change reaches the test, latest first.
\end{itemize}

\section{When to stop}
\label{sec:stop}

A failure ends the run at once. The expensive verdict is \emph{clean}: without any rule, giving it means running
every reached test. A stopping rule saves exactly there, but it has to say how often it is wrong.

\paragraph{The rule.} A stopping rule gives each test in its plan a value $v_1 \ge v_2 \ge \dots$, how much reason is
left to keep going, and stops with ``clean'' at the first value below a threshold $\tau$. For
\strat{adaptive\_bayes} the value before a test is $P(\text{something is broken}) = 1 - P(h = \mathbf{0} \mid
\text{results so far})$.

\paragraph{Setting the threshold.} We set $\tau$ on past regressions, with \emph{split conformal prediction}, a
technique that turns the ranks of past scores into a bound on the error rate. For a past regression $\Delta$, let
$s(\Delta)$ be the value at the first test in the plan that reveals it. With $n$ past regressions and a target miss
rate $\alpha$,
\[
  \tau = s_{(k)}, \qquad k = \lfloor \alpha (n + 1) \rfloor ,
\]
the $k$-th smallest of those values. A new regression is then called clean with probability at most $\alpha$, as
long as past and new changes are \emph{exchangeable}, that is, drawn alike from the same pool of changes.
\begin{itemize}
  \item The bound holds on average, over changes and over the sets of past regressions, and only for changes like
        the past ones.
  \item We set one threshold per group of changes: a skill's steps, text the model reads, an agent's tools and
        skills, the settings, the starting data, the code. A single threshold for all changes can meet the target on
        average while missing most regressions of one group.
\end{itemize}

\paragraph{Properties.} Each test is tagged with the requirements it checks, taken from a catalog of
\emph{properties} (for example, ``a page is never published with broken links''). From the components the change
touches we know which properties it may affect. Before saying ``clean'', the rule runs at least one test, the
cheapest, for every mandatory property the change may affect, and reports which affected properties were checked
and which were not.

\section{Why not a fixed minimal set}
\label{sec:minimal}

The classic alternative is \emph{test-suite minimization}: pick once a small set of tests that still covers the whole
system, and run only that set after every change. With $U_t$ what test $t$ covers (the components it exercises, or the
properties it checks), this is a \emph{set cover} problem:
\[
  \min_{S \subseteq T} |S| \ \ \text{(or } \textstyle\sum_{t \in S} c(t)\text{)} \qquad \text{such that} \qquad
  \bigcup_{t \in S} U_t = \bigcup_{t \in T} U_t ,
\]
usually solved greedily, adding the test that covers the most of what is still uncovered (or the most per token).
Such a set misses regressions for two reasons:
\begin{itemize}
  \item \textbf{Covering is not catching.} A test covers a component when it exercises it, but it catches a broken
        component only if its checks look at exactly the behaviour that broke. A change that skips a permission
        check, for instance, is covered by every test that goes through that code, and caught only by the tests that
        exercise the case where the check should have stopped the action.
  \item \textbf{The set does not know the change.} A fixed set $S$ misses $\Delta$ exactly when $\mathrm{Rev}(\Delta)
        \cap S = \emptyset$. A regression that only one test reveals is caught only if that test happens to be in
        $S$, which was chosen before $\Delta$ existed.
\end{itemize}
Choosing per change avoids both. The replay picks, for each change, the tests through which it can reach the model,
and by Corollary~\ref{cor:safe}, running all of $\Tdiv$ and reading the outcomes of the others from the replay misses
no test that reveals the change. The order then decides what a regression costs, and the threshold decides what a
change that breaks nothing costs, at the miss rate we chose.

\section{Limits}

Besides conditions R and D:
\begin{itemize}
  \item \textbf{Exact matching is conservative.} A request that differs by one character, say a reworded prompt,
        sends the test to the model even if the meaning is the same. When every request carries one long shared
        prompt, editing it reaches every test and the replay saves nothing.
  \item \textbf{The cost bound is an upper bound.} After its divergence point, a run still sends some requests that
        are in $Q_0$, such as the turn of an agent the change does not touch, and the cassette answers those.
  \item \textbf{Clean means clean for this suite.} The replay says which existing tests can see the change. A
        regression that no test checks stays invisible, as with any test selection.
\end{itemize}

\section{Where it lives in the code}

\begin{itemize}
  \item checking one change live: \texttt{traceeval check <change> -{}-history <recording>},
        \texttt{traceeval/experiment/check.py};
  \item the replay: \texttt{traceeval/experiment/replay\_first.py}; the cassette and the clients that answer from
        it: \texttt{traceeval/llm/cache.py};
  \item recording changes on the reached tests only: \texttt{traceeval mutants -{}-reached-only},
        \texttt{traceeval/experiment/runner.py::record\_mutants};
  \item the strategies: \texttt{traceeval/selection/}; the stopping threshold:
        \texttt{traceeval/experiment/stopping.py}.
\end{itemize}

\end{document}
